% Auto-generated by convert2full.sh on 2026-03-10 11:58:03
% Source root: /Users/an/Dropbox/Grants/2026Grants/0430_Margaretha_SEK500k/proposal/tex

% BEGIN FILE: main.tex
\documentclass[11pt,a4paper]{article}

% BEGIN FILE: setup/packages.tex
\usepackage[T1]{fontenc}
\usepackage[utf8]{inputenc}
\usepackage{lmodern}
\usepackage{textcomp}
\usepackage[english]{babel}

\usepackage[
  a4paper,
  left=1.7cm,
  right=1.7cm,
  top=2.0cm,
  bottom=2.0cm,
  headheight=14pt,
  headsep=10pt,
  footskip=16pt
]{geometry}

\usepackage{setspace}
\setstretch{1.0}
\setlength{\parindent}{0pt}
\setlength{\parskip}{5pt}

\usepackage{amsmath,amssymb,amsfonts}
\usepackage{bm}

\usepackage{graphicx}
\usepackage{xcolor}
\usepackage{caption}
\usepackage{subcaption}
\usepackage{booktabs}
\usepackage{array}

\usepackage{tikz}
\usetikzlibrary{arrows.meta,positioning,calc,shapes.geometric}
\usepackage{pgfplots}
\pgfplotsset{compat=1.18}

\usepackage{float}
\usepackage[hidelinks]{hyperref}
\usepackage[nameinlink,noabbrev]{cleveref}
\usepackage{etoolbox}

\usepackage[super,sort&compress]{natbib}
\patchcmd{\thebibliography}{\section*{\refname}}
  {\section*{\centering\refname}}{}{}
\setlength{\bibsep}{1pt}

\usepackage{titlesec}
\titleformat{\section}
  {\normalfont\large\bfseries}
  {\thesection}{0.5em}{}
\titleformat{\subsection}
  {\normalfont\normalsize\bfseries}
  {\thesubsection}{0.5em}{}
\titleformat{\subsubsection}
  {\normalfont\normalsize\itshape}
  {\thesubsubsection}{0.5em}{}
\titlespacing*{\section}{0pt}{8pt plus 2pt minus 1pt}{4pt}
\titlespacing*{\subsection}{0pt}{6pt plus 2pt minus 1pt}{3pt}
\titlespacing*{\subsubsection}{0pt}{5pt plus 1pt minus 1pt}{2pt}
\setcounter{secnumdepth}{3}

\usepackage{enumitem}
\setlist[itemize]{
  topsep=2pt,
  itemsep=1pt,
  parsep=0pt,
  partopsep=0pt,
  leftmargin=1.5em
}
\setlist[enumerate]{
  topsep=2pt,
  itemsep=1pt,
  parsep=0pt,
  partopsep=0pt,
  leftmargin=1.8em
}

\usepackage{fancyhdr}
% END FILE: setup/packages.tex
% BEGIN FILE: setup/metadata.tex
\pagestyle{fancy}
\fancyhf{}

\newcommand{\ApplicantName}{Lijun An}
\newcommand{\ProjectShortTitle}{ProteMapperAtlas}

\fancyhead[L]{\footnotesize \ApplicantName}
\fancyhead[R]{\footnotesize \protect\citep{roberts2019validation}}
\fancyfoot[C]{\footnotesize \thepage}

\renewcommand{\headrulewidth}{0.3pt}
\renewcommand{\footrulewidth}{0pt}

\title{\vspace{-1.8em}\textbf{
ProteMapperAtlas: A Foundation-Style Machine Learning Framework
for Plasma Proteomics in Neurodegenerative Disease
}}
\author{}
\date{}
% END FILE: setup/metadata.tex

\begin{document}

\maketitle
\vspace{-5em}

% BEGIN FILE: sections/01-summary.tex
\section{Summary}

This proposal aims to develop robust and clinically interpretable
artificial intelligence methods for plasma proteomics in
neurodegenerative disease.\citep{roberts2019validation}
The project combines large-scale proteomic profiling, multimodal
phenotyping, and machine learning to improve disease detection,
stratification, and biological understanding.
% END FILE: sections/01-summary.tex
% BEGIN FILE: sections/02-background.tex
\section{Background and State of the Art}

Neurodegenerative diseases impose a growing societal and healthcare
burden. Blood-based proteomics has emerged as a scalable route to
capture systemic and brain-related molecular variation. At the same
time, recent machine learning advances have enabled high-dimensional
pattern extraction, but major limitations remain in transportability,
biological interpretability, and leakage-free evaluation.

A central gap is the lack of principled frameworks that jointly address:
\begin{itemize}
  \item cross-disease modeling,
  \item cross-platform proteomics harmonization,
  \item uncertainty-aware prediction, and
  \item clinically meaningful validation.
\end{itemize}

\subsection{Knowledge gap}

Existing studies often focus on a single disease, a single cohort, or a
single proteomics platform. This limits generalizability and weakens
downstream biological interpretation.

% BEGIN FILE: optional/01-example-figure.tex
% Insert with % Insert with % Insert with % Insert with \input{optional/01-example-figure} inside a subsection.
\begin{figure}[H]
  \centering
  \IfFileExists{figures/example_figure.pdf}{%
    \includegraphics[width=0.82\textwidth]{figures/example_figure.pdf}%
  }{%
    \fbox{%
      \parbox[c][4cm][c]{0.78\textwidth}{\centering
      Placeholder figure: add
      \texttt{figures/example\_figure.pdf}}%
    }%
  }
  \caption{Example inserted figure. Replace the file with your own
  figure in \texttt{figures/}.}
  \label{fig:inserted}
\end{figure}
 inside a subsection.
\begin{figure}[H]
  \centering
  \IfFileExists{figures/example_figure.pdf}{%
    \includegraphics[width=0.82\textwidth]{figures/example_figure.pdf}%
  }{%
    \fbox{%
      \parbox[c][4cm][c]{0.78\textwidth}{\centering
      Placeholder figure: add
      \texttt{figures/example\_figure.pdf}}%
    }%
  }
  \caption{Example inserted figure. Replace the file with your own
  figure in \texttt{figures/}.}
  \label{fig:inserted}
\end{figure}
 inside a subsection.
\begin{figure}[H]
  \centering
  \IfFileExists{figures/example_figure.pdf}{%
    \includegraphics[width=0.82\textwidth]{figures/example_figure.pdf}%
  }{%
    \fbox{%
      \parbox[c][4cm][c]{0.78\textwidth}{\centering
      Placeholder figure: add
      \texttt{figures/example\_figure.pdf}}%
    }%
  }
  \caption{Example inserted figure. Replace the file with your own
  figure in \texttt{figures/}.}
  \label{fig:inserted}
\end{figure}
 inside a subsection.
\begin{figure}[H]
  \centering
  \IfFileExists{../tex/figures/example_figure.pdf}{%
    \includegraphics[width=0.82\textwidth]{../tex/figures/example_figure.pdf}%
  }{%
    \fbox{%
      \parbox[c][4cm][c]{0.78\textwidth}{\centering
      Placeholder figure: add
      \texttt{../tex/figures/example\_figure.pdf}}%
    }%
  }
  \caption{Example inserted figure. Replace the file with your own
  figure in \texttt{../tex/figures/}.}
  \label{fig:inserted}
\end{figure}
% END FILE: optional/01-example-figure.tex
% BEGIN FILE: optional/02-example-tikz.tex
% Insert with % Insert with % Insert with % Insert with \input{optional/02-example-tikz} inside a subsection.
\begin{figure}[H]
\centering
\begin{tikzpicture}[
  node distance=1.8cm and 1.8cm,
  box/.style={
    draw,
    rounded corners,
    align=center,
    minimum width=3.0cm,
    minimum height=0.9cm
  },
  >=Latex
]
\node[box] (data) {Cohorts \\ Proteomics + Clinical Data};
\node[box, right=of data] (model) {Representation \\ Learning Model};
\node[box, right=of model] (outcome)
  {Diagnosis / Prognosis \\ Biological Interpretation};

\draw[->] (data) -- (model);
\draw[->] (model) -- (outcome);
\end{tikzpicture}
\caption{Example schematic drawn directly in LaTeX using TikZ.}
\label{fig:tikz}
\end{figure}
 inside a subsection.
\begin{figure}[H]
\centering
\begin{tikzpicture}[
  node distance=1.8cm and 1.8cm,
  box/.style={
    draw,
    rounded corners,
    align=center,
    minimum width=3.0cm,
    minimum height=0.9cm
  },
  >=Latex
]
\node[box] (data) {Cohorts \\ Proteomics + Clinical Data};
\node[box, right=of data] (model) {Representation \\ Learning Model};
\node[box, right=of model] (outcome)
  {Diagnosis / Prognosis \\ Biological Interpretation};

\draw[->] (data) -- (model);
\draw[->] (model) -- (outcome);
\end{tikzpicture}
\caption{Example schematic drawn directly in LaTeX using TikZ.}
\label{fig:tikz}
\end{figure}
 inside a subsection.
\begin{figure}[H]
\centering
\begin{tikzpicture}[
  node distance=1.8cm and 1.8cm,
  box/.style={
    draw,
    rounded corners,
    align=center,
    minimum width=3.0cm,
    minimum height=0.9cm
  },
  >=Latex
]
\node[box] (data) {Cohorts \\ Proteomics + Clinical Data};
\node[box, right=of data] (model) {Representation \\ Learning Model};
\node[box, right=of model] (outcome)
  {Diagnosis / Prognosis \\ Biological Interpretation};

\draw[->] (data) -- (model);
\draw[->] (model) -- (outcome);
\end{tikzpicture}
\caption{Example schematic drawn directly in LaTeX using TikZ.}
\label{fig:tikz}
\end{figure}
 inside a subsection.
\begin{figure}[H]
\centering
\begin{tikzpicture}[
  node distance=1.8cm and 1.8cm,
  box/.style={
    draw,
    rounded corners,
    align=center,
    minimum width=3.0cm,
    minimum height=0.9cm
  },
  >=Latex
]
\node[box] (data) {Cohorts \\ Proteomics + Clinical Data};
\node[box, right=of data] (model) {Representation \\ Learning Model};
\node[box, right=of model] (outcome)
  {Diagnosis / Prognosis \\ Biological Interpretation};

\draw[->] (data) -- (model);
\draw[->] (model) -- (outcome);
\end{tikzpicture}
\caption{Example schematic drawn directly in LaTeX using TikZ.}
\label{fig:tikz}
\end{figure}
% END FILE: optional/02-example-tikz.tex

\subsection{Why this project is needed}

A unified framework could support both biomarker discovery and disease
mechanism mapping while remaining compatible with real-world clinical
translation.
% END FILE: sections/02-background.tex
% BEGIN FILE: sections/03-preliminary-work.tex
\section{Preliminary Work}

Our group has established expertise in proteomics-based machine
learning, multimodal biomarker modeling, and large-scale
neurodegeneration cohorts. Preliminary analyses show that joint-learning
approaches can improve diagnostic discrimination across multiple
neurodegenerative conditions while identifying biologically plausible
protein signatures.
% END FILE: sections/03-preliminary-work.tex
% BEGIN FILE: sections/04-aim.tex
\section{Aim}

The overall aim is to develop and validate a foundation-style machine
learning framework for plasma proteomics in neurodegenerative disease.

\subsection{Specific Aim 1: Build robust proteomics representation models}

We will train scalable models on large multi-cohort plasma proteomics
datasets to learn transferable protein representations.

\subsection{Specific Aim 2: Improve disease prediction and stratification}

We will evaluate diagnostic and prognostic performance across
Alzheimer’s disease, Parkinson’s disease, frontotemporal dementia,
dementia with Lewy bodies, and related disorders.

\subsection{Specific Aim 3: Link molecular patterns to brain phenotypes}

We will test whether learned protein representations map onto
imaging-derived and cognition-related phenotypes.
% END FILE: sections/04-aim.tex
% BEGIN FILE: sections/05-methods.tex
\section{Research Plan and Methods}

\subsection{Data sources}

The project will use deeply phenotyped cohorts with plasma proteomics,
clinical diagnosis, cognitive testing, and neuroimaging where
available.

\subsection{Proteomics preprocessing}

Protein measurements will undergo rigorous quality control,
normalization, and batch-effect assessment. Missingness patterns will
be characterized explicitly, and model development will be embedded in
leakage-free pipelines.

\subsubsection{Quality control}

Sample- and protein-level exclusions will follow pre-specified
criteria. Outliers, technical artifacts, and panel-specific
inconsistencies will be documented.

\subsubsection{Normalization and harmonization}

We will compare within-platform normalization strategies and
cross-platform alignment approaches where relevant.

\subsection{Model development}

We will evaluate both supervised and representation-learning
approaches. Internal validation will use nested cross-validation or
split-sample designs with strict separation between preprocessing,
feature selection, and evaluation.

\subsection{Statistical analysis}

Performance will be quantified using discrimination, calibration,
uncertainty, and clinical operating characteristics. Subgroup
robustness will be evaluated by age, sex, disease stage, and cohort.
% END FILE: sections/05-methods.tex
% BEGIN FILE: sections/06-expected-results.tex
\section{Expected Results and Scientific Significance}

We expect to generate:
\begin{itemize}
  \item a robust framework for plasma proteomics modeling,
  \item improved disease classification across related
    neurodegenerative disorders,
  \item transferable molecular representations linked to clinically
    relevant phenotypes.
\end{itemize}

Scientifically, the project will clarify how blood-based molecular
patterns reflect disease heterogeneity and brain-related processes.
% END FILE: sections/06-expected-results.tex
% BEGIN FILE: sections/07-innovation.tex
\section{Innovation}

This project is innovative because it integrates:
\begin{itemize}
  \item foundation-style representation learning for proteomics,
  \item cross-disease modeling rather than disease-isolated analysis,
  \item multimodal biological interpretation,
  \item strict leakage-aware benchmarking for trustworthy clinical AI.
\end{itemize}
% END FILE: sections/07-innovation.tex
% BEGIN FILE: sections/08-feasibility.tex
\section{Feasibility}

The proposed work is feasible due to access to large cohorts, strong
computational infrastructure, and established expertise in proteomics
AI and neurodegenerative disease research.
% END FILE: sections/08-feasibility.tex
% BEGIN FILE: sections/09-pitfalls.tex
\section{Potential Pitfalls and Alternative Strategies}

If cross-platform harmonization proves unstable, analyses will
prioritize within-platform foundation models with late-stage alignment.
If certain diagnoses have limited sample size, we will use hierarchical
or transfer-learning strategies.
% END FILE: sections/09-pitfalls.tex
% BEGIN FILE: sections/10-timeline.tex
\section{Timeline}

\begin{itemize}
  \item Year 1: data harmonization, QC, baseline models
  \item Year 2: representation learning, internal validation
  \item Year 3: multimodal interpretation and external validation
  \item Year 4: dissemination, software release, manuscript completion
\end{itemize}
% END FILE: sections/10-timeline.tex
% BEGIN FILE: sections/11-environment.tex
\section{Environment and Resources}

The project will be conducted in a strong interdisciplinary environment
combining clinical neuroscience, biomarker research, and machine
learning expertise. Available infrastructure includes secure data
handling systems, high-performance computing, and access to
longitudinal research cohorts.
% END FILE: sections/11-environment.tex
% BEGIN FILE: sections/12-ethics.tex
\section{Ethics and Data Management}

All analyses will be conducted under approved ethical frameworks and
relevant data-use agreements. Sensitive data will be handled in secure
computational environments. Reproducible pipelines, model cards, and
analysis documentation will be maintained throughout the project.
% END FILE: sections/12-ethics.tex

% BEGIN FILE: sections/13-citations-in-text.tex
\section{References in Text}

Blood-based biomarkers are becoming increasingly important in
neurodegeneration research.\citep{chen2024crossplatform}
Machine learning on high-dimensional omics data requires careful
attention to validation design and transportability.
\citep{chen2024crossplatform}
% END FILE: sections/13-citations-in-text.tex
% BEGIN FILE: references/bibliography.tex
\bibliographystyle{unsrtnat}
\bibliography{../tex/references/references}
% END FILE: references/bibliography.tex

\end{document}
% END FILE: main.tex
